\chapter{Quantum Gravity Candidates}

General relativity is a classical theory: objects are assumed to have definite properties. They are either here or there, in one state or another—not both simultaneously.

In quantum theory, this intuitive picture breaks down. Elementary particles, such as electrons and photons, can exist in superpositions of states—they may be described as being both here and there at once. Only upon measurement do these possibilities resolve into a definite outcome.


\section{Semiclassical Gravity}

One of the earliest attempts to bridge the quantum--classical divide is \emph{semiclassical gravity}. In this approach, matter is treated as fully quantum, while spacetime remains classical. To make the Einstein field equations workable, the operator-valued stress--energy tensor of quantum matter is replaced by its expectation value—the renormalized average of the energy and momentum calculated over the quantum state of the matter fields. This resulting set of ordinary numbers can then be inserted into the equations governing curvature.

Semiclassical gravity is remarkably successful. It accurately describes a wide range of phenomena, from laboratory experiments to astrophysical observations and cosmology. It even predicts striking effects such as Hawking radiation in black holes. Yet its very success also exposes its conceptual limitation: the approach is \emph{ad hoc}.
The theory works well for everything we can observe, but it does not answer any of the deeper question, like what is the physics at the sigularity of a black hole.

\section{Perturbative Quantum Gravity}

A natural next step is \emph{perturbative quantum gravity}, where spacetime is expanded around a simple background—typically flat or slightly curved—and the perturbations are treated as quantum fields. This approach is conceptually straightforward and extends the familiar machinery of quantum field theory to gravity.

However, it quickly runs into a fundamental problem: gravity is nonrenormalizable. Unlike the Standard Model, where quantum infinities can be systematically controlled, perturbative gravity's divergences become increasingly severe at higher energies. To cancel these, we are forced to introduce new counterterms at every order of the calculation. This process never terminates, requiring an infinite number of independent parameters to be measured and fixed. Ultimately, this strips the theory of its predictive power; instead of a concise model, we are left with an infinite list of unknowns. The techniques that work spectacularly well for matter fields simply break down for spacetime itself.

\section{Nonperturbative and Geometric Approaches}

In response to the failure of perturbative quantization, researchers have developed nonperturbative frameworks that do not assume a fixed background geometry. A leading example is Loop Quantum Gravity (LQG), which models spacetime as a discrete combinatorial structure of spin networks. LQG is mathematically rigorous and fully background-independent, offering a conceptually clean quantization of geometry.

However, it seems major obstacles still remain. Deriving a smooth classical spacetime limit is nontrivial, and embedding standard particle physics into the LQG framework remains unresolved. 


\section{Canonical Quantum Gravity: The Wheeler--DeWitt Equation}

A conceptually direct and highly intuitive attempt to quantize gravity is obtained by applying canonical quantization to general relativity.
In this approach, spacetime is decomposed into spatial slices, a bit like a movie of frames, and the Einstein equations are rewritten in Hamiltonian form.

Upon quantization, the classical constraints of general relativity are promoted to operator equations acting on a wavefunctional $\Psi[h_{ij}]$,
which encodes the quantum state of the entire spatial geometry. The central result is the Wheeler--DeWitt equation:
\[
\hat{H}\Psi[h_{ij}] = 0.
\]

We end up to the universe with no time! In standard quantum mechanics, states evolve with respect to time,
but here the wavefunction of the universe appears static. Recovering the notion of time as an emergent or relational concept remains an open problem.



\section{String Theory: A rock too heavy}

In string theory, point particles are replaced by one-dimensional strings. Gravity emerges as one of the vibrational modes of the string, and the framework unifies all fundamental forces in principle. String theory deals with deep mathematical and abstract structures—dualities, extra dimensions, and black-hole entropy counting, to name a few.

However, significant challenges remain. The theory relies on supersymmetry, which predicts that every known particle has a “superpartner” (sparticle) with spin differing by $1/2$. None of these predicted superpartners have been detected experimentally.

The theory also admits an enormous landscape of possible vacuum states—often estimated at around $10^{500}$—raising concerns about predictivity and falsifiability.
These vacua arise from different ways of compactifying extra dimensions, different flux configurations, and different resulting low-energy physical laws, including particle spectra and fundamental constants.
Each vacuum corresponds to a different effective universe. Whether any of these vacua uniquely describes our universe is not known.

Extra spatial dimensions are required, typically assumed to be compactified at extremely small scales, yet they remain empirically undetected. Direct experimental tests of string-scale physics are effectively out of reach.

One might ask what a theory as flexible as string theory ultimately explains. Given enough freedom, one of the $\sim 10^{500}$ vacua may accommodate any observed universe. However, there is no known mechanism to identify the vacuum corresponding to our universe or to derive it uniquely from first principles. In this sense, the theory accommodates reality rather than predicting it.




\section{Emergent and Holographic Approaches}

A more radical class of ideas treats gravity and spacetime as emergent rather than fundamental. This perspective arose from puzzles at the intersection of gravity, thermodynamics, and quantum theory. Black holes behave as thermodynamic objects, possessing entropy proportional to horizon area and emitting thermal radiation. These results suggest a deep link between geometry, information, and statistical mechanics.

The observation that gravitational entropy scales with area rather than volume led to the holographic principle: the idea that the degrees of freedom of a region of spacetime may be encoded on its boundary. Holographic dualities further support this view, showing that spacetime geometry and gravitational dynamics can emerge from nongravitational quantum theories.

The holographic principle argues that a three-dimensional universe can be described by a two-dimensional theory ($N \rightarrow N-1$). This is actually quite surprising result.
Everything that happens in the universe (whether it had four dimensions or eleveven, can be described by its $N-1$ dimensional surface. 


\section{The Simulation Hypothesis: Reality as Software}

During recent years so-called Simulation Hypothesis has become increasingly popular. If $3D$ can map to $2D$, why stop there? In software architecture, any $N$-dimensional space is ultimately stored as a one-dimensional bitstring ($N \rightarrow 1$). A sequence of bits has no intrinsic geometry, it is just a set of information.

If the universe is fundamentally informational, it is tempting to conclude that we are merely a program running on some higher-order "hardware." In this view, the strange "quantization" of our world is simply the resolution of the grid, and the speed of light is the clock-speed of the processor.

However, the simulation hypothesis feels like a philosophical "shell game." It merely translates the mystery of existence by one level: if we are a simulation, who simulated the simulators? Furthermore, it ignores the staggering Information Cost of reality.

Consider the entropy of a single human being. To simulate even a single strand of DNA with perfect fidelity requires tracking billions of quantum interactions. To harvest enough information from a "parent universe" to simulate an entire "child universe" would require a massive thermodynamic overhead. Who would pay the electricity bill for that project?

\textit{The simulation hypothesis might explain why General Relativity and Quantum Mechanics appear so difficult to unify. Anyone familiar with software engineering is aware of legacy code. The fundamental incompatibility between the smooth, geometric curves of GR and the discrete, probabilistic jumps of QM might not be a profound mystery of nature, but simply poor design. In this light, the universe is a patchwork of modules written by different architects, at different times, with different goals.}


